% GENERATED FILE. Edit canonical lessons, then run npm run generate:books.
% canonical-source-hash: fnv1a64:e5f7d86acf0edfec
% canonical-lessons: KA-C254-telu2, KA-C254-asvastha, KA-C254-cupada, KA-C254-kandu, KA-C254-gulabi

\chapter{Thin, Unwell, Sharp, Brown, Pink}
\label{ch:ka-thin-unwell-sharp-brown-pink}
\hlchaptermodality{254}

\begin{chapteropening}
\textbf{By the end of this chapter:} \emph{I can say thin, unwell, sharp, brown and pink.}

Complete the last lesson of chapter 254: I can say thin, unwell, sharp, brown and pink.
\end{chapteropening}

\section[\texorpdfstring{\kn{ತೆಳು} (teḷu)}{ತೆಳು (teḷu)}]{\kn{ತೆಳು} (teḷu) — thin}

\label{lesson:KA-C254-telu2}



\begin{quote}
Before the new one: say the Kannada for life, then the Kannada for a rule.
\end{quote}

\subsection*{You'll want to know: \kn{ತೆಳು}}
\textbf{\kn{ತೆಳು}} — \emph{teḷu} — \textquotedblleft{}thin\textquotedblright{}.

\begin{grammarlens}[title={What you've built}]
One more word for this chapter.
\end{grammarlens}

\subsection*{Your turn}
\begin{itemize}
\raggedright
  \item \emph{Say it:} \emph{teḷu}
  \item \emph{Say it:} \emph{teḷu}, once more
  \item \emph{Say it:} say \emph{teḷu}
  \item \emph{Recall:} say the Kannada for life, then the Kannada for a rule, then say \emph{teḷu} again
\end{itemize}

\begin{tcolorbox}[breakable,colback=teal!4,colframe=teal!35!black,title={Before you move on}]
What does \kn{ತೆಳು} mean? (\textquotedblleft{}Thin\textquotedblright{}.) Say it once more.
\end{tcolorbox}
\section[\texorpdfstring{\kn{ಅಸ್ವಸ್ಥ} (asvastha)}{ಅಸ್ವಸ್ಥ (asvastha)}]{\kn{ಅಸ್ವಸ್ಥ} (asvastha) — unwell, sick}

\label{lesson:KA-C254-asvastha}



\begin{quote}
Before the new one: say the Kannada for a rule, then the Kannada for thin.
\end{quote}

\subsection*{You'll want to know: \kn{ಅಸ್ವಸ್ಥ}}
\textbf{\kn{ಅಸ್ವಸ್ಥ}} — \emph{asvastha} — \textquotedblleft{}unwell, sick\textquotedblright{}.

\begin{grammarlens}[title={What you've built}]
One more word for this chapter.
\end{grammarlens}

\subsection*{Your turn}
\begin{itemize}
\raggedright
  \item \emph{Say it:} \emph{asvastha}
  \item \emph{Say it:} \emph{asvastha}, once more
  \item \emph{Say it:} say \emph{asvastha}
  \item \emph{Recall:} say the Kannada for a rule, then the Kannada for thin, then say \emph{asvastha} again
\end{itemize}

\begin{tcolorbox}[breakable,colback=teal!4,colframe=teal!35!black,title={Before you move on}]
What does \kn{ಅಸ್ವಸ್ಥ} mean? (\textquotedblleft{}Unwell, sick\textquotedblright{}.) Say it once more.
\end{tcolorbox}
\section[\texorpdfstring{\kn{ಚೂಪಾದ} (cūpāda)}{ಚೂಪಾದ (cūpāda)}]{\kn{ಚೂಪಾದ} (cūpāda) — sharp}

\label{lesson:KA-C254-cupada}



\begin{quote}
Before the new one: say the Kannada for thin, then the Kannada for unwell.
\end{quote}

\subsection*{You'll want to know: \kn{ಚೂಪಾದ}}
\textbf{\kn{ಚೂಪಾದ}} — \emph{cūpāda} — \textquotedblleft{}sharp\textquotedblright{}.

\begin{grammarlens}[title={What you've built}]
One more word for this chapter.
\end{grammarlens}

\subsection*{Your turn}
\begin{itemize}
\raggedright
  \item \emph{Say it:} \emph{cūpāda}
  \item \emph{Say it:} \emph{cūpāda}, once more
  \item \emph{Say it:} say \emph{cūpāda}
  \item \emph{Recall:} say the Kannada for thin, then the Kannada for unwell, then say \emph{cūpāda} again
\end{itemize}

\begin{tcolorbox}[breakable,colback=teal!4,colframe=teal!35!black,title={Before you move on}]
What does \kn{ಚೂಪಾದ} mean? (\textquotedblleft{}Sharp\textquotedblright{}.) Say it once more.
\end{tcolorbox}
\section[\texorpdfstring{\kn{ಕಂದು} (kandu)}{ಕಂದು (kandu)}]{\kn{ಕಂದು} (kandu) — brown}

\label{lesson:KA-C254-kandu}



\begin{quote}
Before the new one: say the Kannada for unwell, then the Kannada for sharp.
\end{quote}

\subsection*{You'll want to know: \kn{ಕಂದು}}
\textbf{\kn{ಕಂದು}} — \emph{kandu} — \textquotedblleft{}brown\textquotedblright{}.

\begin{grammarlens}[title={What you've built}]
One more word for this chapter.
\end{grammarlens}

\subsection*{Your turn}
\begin{itemize}
\raggedright
  \item \emph{Say it:} \emph{kandu}
  \item \emph{Say it:} \emph{kandu}, once more
  \item \emph{Say it:} say \emph{kandu}
  \item \emph{Recall:} say the Kannada for unwell, then the Kannada for sharp, then say \emph{kandu} again
\end{itemize}

\begin{tcolorbox}[breakable,colback=teal!4,colframe=teal!35!black,title={Before you move on}]
What does \kn{ಕಂದು} mean? (\textquotedblleft{}Brown\textquotedblright{}.) Say it once more.
\end{tcolorbox}
\section[\texorpdfstring{\kn{ಗುಲಾಬಿ} (gulābi)}{ಗುಲಾಬಿ (gulābi)}]{\kn{ಗುಲಾಬಿ} (gulābi) — pink}

\label{lesson:KA-C254-gulabi}



\begin{quote}
Before the new one: say the Kannada for sharp, then the Kannada for brown.
\end{quote}

\subsection*{You'll want to know: \kn{ಗುಲಾಬಿ}}
\textbf{\kn{ಗುಲಾಬಿ}} — \emph{gulābi} — \textquotedblleft{}pink\textquotedblright{}.

\begin{grammarlens}[title={What you've built}]
One more word for this chapter.
\end{grammarlens}

\subsection*{Your turn}
\begin{itemize}
\raggedright
  \item \emph{Say it:} \emph{gulābi}
  \item \emph{Say it:} \emph{gulābi}, once more
  \item \emph{Say it:} say \emph{gulābi}
  \item \emph{Recall:} say the Kannada for sharp, then the Kannada for brown, then say \emph{gulābi} again
\end{itemize}

\begin{tcolorbox}[breakable,colback=teal!4,colframe=teal!35!black,title={Before you move on}]
What does \kn{ಗುಲಾಬಿ} mean? (\textquotedblleft{}Pink\textquotedblright{}.) Say it once more.
\end{tcolorbox}
